\section{Free exterior matching}\label{sec:free}

The profile construction and the spectral count both use a free radial
equation in dimension twelve. We study it first as an independent matching
problem. For a real parameter $b$, a spectral parameter $\lambda\in\mathbb C$,
an angular degree $\ell\in\{0,1,2,\ldots\}$, and a channel index $h=\pm1$,
the equation is
\begin{equation}\label{free:radial-equation}
 v_h''+\left(\frac{11}{r}+\frac{hi r}{2}\right)v_h'
 +\left(b+hi\lambda-\frac{\ell(\ell+10)}{r^2}\right)v_h=0.
\end{equation}
Here $\ell(\ell+10)$ is the eigenvalue of the negative spherical Laplacian
on degree-$\ell$ harmonics on $\mathbb S^{11}$. The substitution
$v_h=r^\ell H(q,6+\ell,x)$, with
$q=\lambda+\ell/2-hib$ and $x=-hi r^2/4$, gives the confluent
hypergeometric equation. Its slow solution is
\[
 r^\ell H(\lambda+\ell/2-hi b,6+\ell,-hi r^2/4),
\]
where Lemma~\ref{free:slow} defines $H$ by its integral and its
nonoscillatory asymptotic expansion. Primes on $v_h$ denote radial
derivatives; primes on $H$ denote derivatives in its last argument.
We need two kinds of boundary data
from this solution. At $\lambda=\ell=0$ in channel $h=1$, the quotient
$j=-H'/H$ vanishes exactly when the radial derivative vanishes; this
matches the exterior to a constant core. For general $\lambda,\ell$,
the two value-and-derivative columns in Proposition~\ref{free:main}
define the matching determinant whose zeros we count.

Section~\ref{sec:profiles} derives the similarity equation and constructs
the finite-power profiles. In those variables $b$ is the logarithmic phase
frequency and a spectral mode has time factor $e^{\lambda s}$, where
$s$ is logarithmic similarity time.
Lemma~\ref{spec:interface} will identify the determinant studied here
with the limiting constrained spectral problem. The calculations in this
section concern its free exterior, where the power parameter is $a=0$
and the pressure is zero.

All terminating decimals in this section denote exact rational numbers.  Set
\begin{equation}\label{free:disk}
 \begin{gathered}
 b_*=.33477606871236,\qquad Z_*=2.70506819293654,\qquad
 \delta=10^{-8},\\
 \mathcal D=\{(b,Z)\in\mathbb R^2:(b-b_*)^2+(Z-Z_*)^2\le\delta^2\}.
 \end{gathered}
\end{equation}
The radius of the limiting core is $R=2\sqrt Z$.
We prove the following statement, including its finite
arithmetic ingredients.  The definition of $H$ follows immediately below.
The degree on $\mathcal D$ means Brouwer degree on its interior.

\begin{proposition}[Free matching]\label{free:main}
For every $(b,Z)\in\mathcal D$, the following conclusions hold.
\begin{enumerate}
\item The quotient
\[
 j(b,Z)=-\frac{H'(-ib,6,-iZ)}{H(-ib,6,-iZ)}
\]
is well defined and satisfies $|j(b,Z)|<3\cdot10^{-9}$.  It is nonzero on
$\partial\mathcal D$, and, identifying $\mathbb C$ with $\mathbb R^2$ in the
order real part, imaginary part,
\[
 \deg(j,\mathcal D,0)=-1.
\]
\item For every real $z\ge Z$, $H(-ib,6,-iz)\ne0$ and
\begin{equation}\label{free:exterior-bounds}
 \operatorname{Re}j(b,z)<.043.
\end{equation}
For $z\ge3$ one also has $\operatorname{Im}j(b,z)>0$.
\item For $\ell\in\{0,1,2,\ldots\}$ and $h\in\{1,-1\}$ define
\begin{equation}\label{free:jet}
 q_h=\lambda+\ell/2-hi b,\qquad
 J_h(\lambda)=
 \begin{pmatrix}H(q_h,6+\ell,-hiZ)\\-H'(q_h,6+\ell,-hiZ)\end{pmatrix}.
\end{equation}
The columns $J_h$ are holomorphic and nonzero on a neighborhood of
$\operatorname{Re}\lambda\ge-1/32$.  The determinant
$\Delta_\ell(\lambda;b,Z)=\det(J_+(\lambda),J_-(\lambda))$ has no zero on
$\operatorname{Re}\lambda=-1/32$.  In the closed half-plane
$\operatorname{Re}\lambda\ge-1/32$ its numbers of zeros, counted with
holomorphic multiplicity, are
\[
 \begin{array}{c|ccc}
 \ell&0&1&\ell\ge2\\ \hline
 \#\{\Delta_\ell=0\}&2&1&0.
 \end{array}
\]
\end{enumerate}
The locations of the symmetry zeros are identified later at parameters
where $j(b,Z)=0$; the count in this proposition holds on the entire disk.
\end{proposition}

\subsection{The slow solution and the Laguerre cone}

The function $H(q,m,x)$ below is Tricomi's confluent hypergeometric
function $U(q,m,x)$; see \cite[Equations~13.4.4 and~13.3.22]{DLMF}.
We prove the regularized representation and the uniform symbol bounds
needed on the imaginary rays.

\begin{lemma}[Regularized slow solution]\label{free:slow}
For $\operatorname{Re}q>-1$, $\operatorname{Re}x\ge0$, $x\ne0$, define, with
principal powers,
\begin{equation}\label{free:H-definition}
 H(q,m,x)=x^{-q}\left\{1+\frac1{\Gamma(q)}
 \int_0^\infty e^{-u}u^{q-1}
 \left[(1+u/x)^{m-1-q}-1\right]\,du\right\}.
\end{equation}
Here $m$ may in particular be any positive integer used below.  This
definition is holomorphic in $q$, including $q=0$, and in $x$ on the
interior of the indicated region, with holomorphic continuation across
each nonzero point of its imaginary boundary.  Primes on $H$ denote
derivatives in its last argument.  It satisfies
\begin{equation}\label{free:H-equation}
 xH''+(m-x)H'-qH=0,\qquad
 H'(q,m,x)=-qH(q+1,m+1,x),\qquad H(0,m,x)=1.
\end{equation}
For every $N\ge1$ it has the expansion
\begin{equation}\label{free:H-expansion}
 H(q,m,x)=x^{-q}\left\{
 \sum_{j=0}^{N-1}\frac{(-1)^j(q)_j(1+q-m)_j}{j!x^j}
 +O(|x|^{-N})\right\}.
\end{equation}
Here $(q)_0=1$ and $(q)_j=\prod_{n=0}^{j-1}(q+n)$ for $j\ge1$.
The expansion is uniform for $\operatorname{Re}x\ge0$, $|x|\ge1$, and compact parameter
sets.  The same remainder estimate holds after any fixed number of
$x\partial_x$ derivatives of the bracket.  The symbol expansion specifies
the slow solution, without an oscillatory component on the imaginary
rays.
\end{lemma}

\begin{proof}
The bracket in the integrand of \eqref{free:H-definition} is $O(u)$ at
zero, locally uniformly in the parameters; at infinity it and its
parameter derivatives grow at most polynomially, with additional powers
of $\log u$.  Thus the regularized integral is locally holomorphic for
$\operatorname{Re}q>-1$.  The reciprocal gamma function vanishes at zero,
which proves the last assertion in \eqref{free:H-equation}.

For $\operatorname{Re}q>0$ the expression in braces can be combined into
the usual integral.  On $x>0$, substitution $u=xv$ gives
\[
 H(q,m,x)=\frac1{\Gamma(q)}\int_0^\infty
 e^{-xv}v^{q-1}(1+v)^{m-1-q}\,dv.
\]
Differentiation and integration by parts prove both the equation and the
derivative identity; holomorphic continuation proves them on the stated
domain.  The branches in the regularized integral are regular on a
neighborhood of each nonzero imaginary $x$, so the boundary values used
in the matching problem are genuine holomorphic values.

To obtain \eqref{free:H-expansion}, split the integral at $u=|x|/2$.
Taylor's formula applies uniformly to $(1+u/x)^{m-1-q}$ on the first
piece.  Its $j$th coefficient integrates to
$(-1)^j(q)_j(1+q-m)_j/j!$.  On the second piece the exponential weight
bounds the integral by any prescribed inverse power of $|x|$; the same
argument applies after symbol differentiation.  This proves the full
assertion about remainders.  For completeness, the other local solution
at infinity obtained by reduction of order has the form
$e^xx^{q-m}$ times an inverse-power expansion.  On either imaginary ray
its repeated symbol derivatives acquire arbitrarily high powers of
$|x|$.  It therefore cannot be added to \eqref{free:H-expansion} while
preserving all its symbol estimates.  This also explains the slow
normalization on the rays.
\end{proof}

Fix $h,\ell,b,Z$ and write
\[
 M=\ell+5,\quad s=hiZ,\quad q=\lambda+\ell/2-hi b,
 \quad G(t)=H(q,M+1,t-s)\quad(t\ge0).
\]
All parameter analyticity statements below concern $\lambda$; estimates
are locally uniform in the other displayed parameters.

\begin{lemma}[Laguerre recurrence]\label{free:laguerre}
Let $L_n(t)=e^t\partial_t^n(e^{-t}t^n)/n!$ be the Laguerre polynomials,
orthonormal for $e^{-t}\,dt$ on $[0,\infty)$.  Write
\[
 G=\sum_{n\ge0}g_nL_n,\qquad
 B_n=\sum_{k\ge n}g_k,\qquad C_n=\sum_{k\ge n}(k-n)g_k.
\]
These series and their parameter derivatives converge rapidly.  With
$t_n=q+n$ and $d_n=t_n-M$, one has
\begin{align}
 &(B_0,C_0)=(G(0),-G'(0)),\label{free:initial-tail}\\
 &t_ng_n=-MB_{n+1}-sC_n,\qquad
 d_nB_{n+1}=t_nB_n+sC_n,\qquad
 C_{n+1}=C_n-B_{n+1}.\label{free:recurrence}
\end{align}
The real quadratic form
\begin{equation}\label{free:cone-form}
 p_n=\frac M2|B_n|^2+\operatorname{Re}(s\overline{B_n}C_n)
\end{equation}
satisfies the exact identity
\begin{equation}\label{free:cone-increment}
 p_{n+1}-p_n=
 (\operatorname{Re}\lambda+n-5/2)|g_n|^2.
\end{equation}
\end{lemma}

\begin{proof}
Integration by parts gives orthonormality and
$[t\partial_t^2+(1-t)\partial_t]L_n=-nL_n$.
Polynomials are complete: a function orthogonal to all powers has a
Laplace transform whose derivatives vanish at an interior point of its
domain of holomorphy, hence its Laplace transform, and then the function,
vanishes.  The function $G$ and all its derivatives are smooth at zero
and of polynomial growth at infinity by Lemma~\ref{free:slow}.
Repeated integration by parts with this Laguerre operator shows that
$g_n$, and the Laguerre coefficients of each derivative of $G$, decrease
faster than every inverse power of $n$.  The boundary terms vanish
because the coefficient of the derivative at zero is $t$, and the
weight at infinity is exponential.  These arguments are uniform on
compact parameter sets, also after holomorphic differentiation.

The generating function gives, initially for $0\le x<1$,
\[
 \sum_{n\ge0}g_nx^n=
 \int_0^\infty e^{-u}G((1-x)u)\,du.
\]
Taking its value and first derivative as $x\uparrow1$ proves
\eqref{free:initial-tail}.  Since
$L_k'=-\sum_{n<k}L_n$, the coefficients of $G'$ are $-B_{n+1}$,
whereas those of $G''-G'$ are $C_n$.  Rewriting the differential equation
as
\[
 [t\partial_t^2+(1-t)\partial_t]G+MG'-s(G''-G')-qG=0
\]
proves the first recurrence.  The other two follow from
$g_n=B_n-B_{n+1}$ and the definition of $C_n$.

Finally insert $B_{n+1}=B_n-g_n$ and
$C_{n+1}=C_n-B_{n+1}$ into \eqref{free:cone-form}.  Since $s$ is purely
imaginary, the cross term $\operatorname{Re}(s|B_{n+1}|^2)$ vanishes.
Use $t_ng_n=-MB_{n+1}-sC_n$ to simplify the remaining terms.  The
coefficient of $|g_n|^2$ becomes
$\operatorname{Re}t_n-M/2=\operatorname{Re}\lambda+n-5/2$, which is
\eqref{free:cone-increment}.
\end{proof}

\begin{lemma}[Strict normalized cone]\label{free:cone}
Suppose $\operatorname{Re}\lambda\ge-1/32$ and $K\ge3$.  For $n\ge K$
the normalized coefficients $\widehat g_n=g_n/(q)_K$, and their sums
$\widehat B_K=B_K/(q)_K$, $\widehat C_K=C_K/(q)_K$, extend
holomorphically across $q=0$.  They satisfy
\begin{equation}\label{free:strict-cone}
 \frac M2|\widehat B_K|^2+
 \operatorname{Re}(s\overline{\widehat B_K}\widehat C_K)
 =-\sum_{n\ge K}(\operatorname{Re}\lambda+n-5/2)
 |\widehat g_n|^2<0.
\end{equation}
Consequently $\widehat C_K\ne0$, and
\begin{equation}\label{free:tail-disk}
 R_h=\widehat B_K/\widehat C_K\quad\hbox{is holomorphic},\qquad
 |R_h+s/M|<Z/M.
\end{equation}
\end{lemma}

\begin{proof}
The only zero of $(q)_K$ in the specified half-plane is $q=0$, and it is
simple.  It can occur only for $\ell=0$.  At this point $G=1$, so every
$g_n$ with $n\ge1$ vanishes.  Holomorphic division therefore removes
the apparent singularity for $n\ge K$.  Local Cauchy estimates preserve
the rapid decay established above, justifying division of the sums and
continuation of their identities.

The normalized tail cannot vanish identically in its sequence index.
Indeed repeated use of \eqref{free:H-equation} gives
\begin{equation}\label{free:normalized-derivative}
 \frac{\partial_t^K G}{(q)_K}
 =(-1)^KH(q+K,M+1+K,t-s).
\end{equation}
The right side is nonzero by its leading asymptotic coefficient, also at
$q=0$.  If all normalized coefficients with $n\ge K$ vanished, the
function $[G-\sum_{n<K}g_nL_n]/(q)_K$, with its removable value at zero,
would vanish.  Its $K$th derivative would be zero, contradicting
\eqref{free:normalized-derivative}.

Sum \eqref{free:cone-increment} from $K$ to infinity and divide by
$|(q)_K|^2$, taking the removable value when necessary.  Rapid decay
makes the limiting form zero.  Each weight in the resulting sum is
positive because $K\ge3$ and $\operatorname{Re}\lambda\ge-1/32$.
The preceding nontriviality gives strictness in
\eqref{free:strict-cone}.  If $\widehat C_K$ were zero the left side
would be nonnegative.  Dividing by its nonzero squared modulus and
completing the square proves \eqref{free:tail-disk}.
\end{proof}

\begin{lemma}[Analytic matching columns]\label{free:columns}
Define
\begin{equation}\label{free:backward}
 F_n=\begin{pmatrix}d_n-s&-s\\t_n&t_n\end{pmatrix},\qquad
 P_K=F_0F_1\cdots F_{K-1}.
\end{equation}
Then, including parameters where an individual $t_n$ or $d_n$ vanishes,
\begin{equation}\label{free:column-identity}
 P_K\begin{pmatrix}R_h\\1\end{pmatrix}
 =\frac1{\widehat C_K}
 \begin{pmatrix}H(q,M+1,-s)\\-H'(q,M+1,-s)\end{pmatrix}.
\end{equation}
The scalar on the right is holomorphic and nonzero.  The column is
nonzero.
\end{lemma}

\begin{proof}
The recurrences imply the division-free identity
$F_n(B_{n+1},C_{n+1})^t=t_n(B_n,C_n)^t$.  Multiplying gives
$P_K(B_K,C_K)^t=(q)_K(B_0,C_0)^t$.  Divide by $(q)_K$ and use
Lemma~\ref{free:cone}; analytic continuation gives the identity at the
removable point.  A simultaneous zero of $H$ and $H'$ at $x=-s\ne0$
would, by uniqueness for its second order equation, make $H$ identically
zero, contrary to \eqref{free:H-expansion}.
\end{proof}

\subsection{Exclusion of high angular degrees}

\begin{lemma}\label{free:high-angular}
For $(b,Z)\in\mathcal D$, $\ell\ge4$, and
$\operatorname{Re}\lambda\ge-1/32$, one has $\Delta_\ell(\lambda;b,Z)\ne0$.
\end{lemma}

\begin{proof}
Suppose the determinant vanished.  Since the two columns are nonzero,
multiply their functions by nonzero constants to arrange
$H_+(-iZ)=H_-(iZ)$ and $H_+'(-iZ)=H_-'(iZ)$.
Here $H_h(x)$ denotes the resulting multiple of
$H(\lambda+\ell/2-hi b,M+1,x)$ and $M=\ell+5\ge9$.
The substitution $U_h=e^{-x/2}x^{M/2}H_h$ gives
\begin{equation}\label{free:gauged}
 -(x\partial_xU_h)_x+\left(\frac x4+\frac{M^2}{4x}\right)U_h
 =(3-\lambda+hi b)U_h.
\end{equation}

Integrate this equation along the conjugate paths
\[
 \begin{gathered}
 x_h(\eta)=u-hi v,\quad x_h(0)=-hiZ,\quad
 g_h=x_h'=e^{-hi\vartheta},\\
 \vartheta=\min(2\eta,\vartheta_0),\quad
 \cos\vartheta_0=\frac{99}{101},\quad
 \sin\vartheta_0=\frac{20}{101}.
 \end{gathered}
\]
Write $\eta_0=\vartheta_0/2$ and
\[
 x_h/g_h^2=L+hiN,\qquad
 L=u\cos(2\vartheta)+v\sin(2\vartheta),\quad
 N=u\sin(2\vartheta)-v\cos(2\vartheta).
\]
In particular $L>0$ for $\eta>0$ and
\[
 L'=\cos\vartheta-2\vartheta'N,\qquad
 N'=\sin\vartheta+2\vartheta'L.
\]
The real part of $x$ tends to $+\infty$ linearly, so the exponential
factor in $U_h$ justifies every boundary limit and integration below.

For a bounded real function $\alpha$, smooth on each side of $\eta_0$,
put $B_c=\alpha L+\vartheta'N/2$.  Pair
\eqref{free:gauged} with $\overline U_h\,d\eta$ and take real parts.
The kinetic expression after integration by parts is
\[
 L|U_h'|^2+\operatorname{Re}
 \big(hi\vartheta'(L+hiN)\overline U_hU_h'\big).
\]
Completing the square writes this as
$L|U_h'+(\alpha-hi\vartheta'/2)U_h|^2
-B_c(|U_h|^2)'-L(\alpha^2+\vartheta'^2/4)|U_h|^2$.
The real lower endpoint terms before this completion sum to
$\sum_h(M/2)|U_h(0)|^2$.  To see the cancellation, expand
$x\overline U_h\partial_xU_h$ using $U_h=e^{-x/2}x^{M/2}H_h$:
the terms containing $x\overline H_hH_h'$ cancel between the two
matched jets, and the real contribution from the gauge is $M/2$.
Since $B_c(0)=-Z$, integration of the displayed derivative, separately
on the two path pieces, now yields
\begin{align}\label{free:multiplier-identity}
 0=\sum_{h=\pm1}\Bigg\{&
 (M/2-Z)|U_h(0)|^2+
 [B_c(\eta_0+)-B_c(\eta_0-)]|U_h(\eta_0)|^2\notag\\
 &+\int_0^\infty
 \left[L|U_h'+(\alpha-hi\vartheta'/2)U_h|^2
       +V_c|U_h|^2\right]\,d\eta\Bigg\},
\end{align}
where, with $\sigma=\operatorname{Re}\lambda$,
\begin{equation}\label{free:multiplier-potential}
 V_c=\sigma-3+\frac u4+\frac{M^2u}{4(u^2+v^2)}
 +\frac{\vartheta'\sin\vartheta}{2}
 +\frac34L\vartheta'^2+(\alpha L)'-L\alpha^2.
\end{equation}
There is no missing corner term: it is precisely the jump displayed in
\eqref{free:multiplier-identity}.

On the arc take $\alpha=.3$.  Here
$u=\sin\vartheta/2$, $v=Z+(1-\cos\vartheta)/2$, and
\[
 u\le.1,\qquad 2.70<v<2.72,\qquad
 N\le.04-2.70\cdot.92=-2.444,
 \qquad L'>10.7.
\]
It follows directly from \eqref{free:multiplier-potential}, dropping
its nonnegative $u$ terms and sine term, that
$V_c>-3-1/32+.3(10.7)+(3-.09)L>0$.
At the join $1.13<L(\eta_0)<1.19$.  On the straight piece, with
$t=\eta-\eta_0$, choose
\[
 \alpha L=-1.5+\frac{5t}{1+1.5t}.
\]
The jump is $-1.5-.3L(\eta_0)-N(\eta_0)>0$.
The geometry of the ray gives
\[
 u_0=.099+.98t\le u\le.100+.981t,\qquad
 v\le2.72+.20t,\qquad L\ge L_0=1.13+.98t.
\]
Set $d_0=(.100+.981t)^2+(2.72+.20t)^2$ and $q_0=1+1.5t$.
Since $M^2\ge81$, $V_cL_0d_0q_0^2$ is bounded below by
\begin{align}\label{free:high-polynomial}
 P(t)={}&(-3-1/32+u_0/4)L_0d_0q_0^2
 +\frac{81}{4}u_0L_0q_0^2+5L_0d_0
 -(-1.5+2.75t)^2d_0\notag\\
 >&2+t(33-72.2t+44.8t^2)
       +t^4(17.9-4.61t+.54t^2)>0.
\end{align}
The first strict comparison is coefficientwise exact arithmetic; the
last two quadratic discriminants are $-17519/25$ and
$-174119/10000$.  Appendix~\ref{cert:high-angular} provides the full
polynomial recipe and verifier.  Thus all terms of
\eqref{free:multiplier-identity} are nonnegative, the bulk potential is
strictly positive, and $M/2>Z$.  It follows that $U_h=0$, a
contradiction.
\end{proof}

\subsection{Boundary homotopy and the low angular counts}

We use $K=8$ for $0\le\ell\le3$.  For $0\le\rho\le1$ replace the tail
ratio in \eqref{free:column-identity} by $\rho R_h$.  The closure of
the disk in \eqref{free:tail-disk} is convex and contains zero, so every terminal
pair $(\rho R_h,1)$ satisfies $p_K\le0$.  The resulting columns are
holomorphic functions of $\lambda$, continuous in $\rho$.  We must
exclude determinant zeros on a common counting contour.

\begin{lemma}[Counting boundary]\label{free:boundary}
There is a rectangle with left edge $\operatorname{Re}\lambda=-1/32$
and other edges sufficiently far to the right and above and below,
uniformly for $(b,Z)\in\mathcal D$ and $0\le\rho\le1$, on whose
boundary the homotoped determinant is nonzero.  The same exclusions
hold on every larger such rectangle.
\end{lemma}

\begin{proof}
Choose the edges so that $t_n,d_n\ne0$ for $n<8$ on them.  If the
columns became dependent at such a point, their nonzero common initial
jet $(B_0,C_0)$ would have two forward images satisfying $p_{8,h}\le0$.
Indeed the backward matrices are invertible there and their inverses
are exactly the forward recurrences, up to nonzero scalar factors.
For the same initial jet in both channels,
\begin{equation}\label{free:summed-increment}
 p_{8,+}+p_{8,-}=M|B_0|^2+
 \sum_{h=\pm1}\sum_{n=0}^7
 (\sigma+n-5/2)|g_{n,h}|^2.
\end{equation}
For $\sigma>5/2$ all increment weights are positive.  If $B_0\ne0$
the first term is positive; if $B_0=0$ and $C_0\ne0$, then
$g_{0,h}=-s_hC_0/d_{0,h}\ne0$.  This excludes the right edge,
which can be taken beyond every possible zero of its divisors.

For the horizontal edges let $v=\operatorname{Im}\lambda$,
$C_0=vc$, and normalize $|B_0|^2+|c|^2=1$.
Uniformly for $\sigma$ in a bounded interval and $n<8$,
\[
 g_{0,h}=-(MB_0+s_hC_0)/d_{0,h}=-hZc+o(1),\qquad
 g_{n+1,h}=\frac{t_{n,h}g_{n,h}+s_hB_{n+1,h}}{d_{n+1,h}}
             =g_{n,h}+o(1).
\]
Consequently \eqref{free:summed-increment} equals
\[
 M|B_0|^2+2(8\sigma+8)Z^2|c|^2+o(1)>0
\]
for large $|v|$, since $8\sigma+8\ge31/4$.
This proves the two horizontal exclusions uniformly.

For the left edge set $s=iZ$ and form the forward matrix without its
divisors:
\begin{equation}\label{free:forward-matrix}
 T=T_{K-1}\cdots T_0,\qquad
 T_n=\begin{pmatrix}t_n&s\\-t_n&d_n-s\end{pmatrix},\qquad K=8.
\end{equation}
Thus the actual forward jet is $T(B_0,C_0)^t/\prod_{n<K}d_n$.
Define
\begin{equation}\label{free:boundary-form}
 \mathsf H(v)=T^*
 \begin{pmatrix}M&s\\-s&0\end{pmatrix}T,
 \qquad \mathsf S(v)=\mathsf H(v)+\overline{\mathsf H(-v)}.
\end{equation}
The second summand is the form for $h=-1$.  Each summand evaluates to
$2|\prod d_{n,h}|^2p_{K,h}$ on the common initial pair.  We shall prove
$\mathsf S(v)>0$ for every real $v$, excluding the two simultaneous
nonpositive values.

First, its determinant is a real even polynomial of degree at most
$4K-2=30$ for every real $(b,Z)$ in the parameter box.  Hermitian
symmetry gives real diagonal entries, and the reflection in
\eqref{free:boundary-form} makes the diagonal entries even and the
off-diagonal entry satisfy $C(-v)=\overline{C(v)}$.  To check the degree
bound structurally, hold the initial jet fixed and let $|v|\to\infty$.
The forward recurrences give $g_{n,h}=O((|B_0|+|C_0|)/|v|)$ for
$n<K$.  Also $|\prod d_{n,h}|^2=|v|^{2K}(1+O(1/|v|))$.
Using \eqref{free:cone-increment}, each form divided by $|v|^{2K}$
has its initial cone form as leading term and an $O(v^{-2})$
increment.  The opposite initial cross terms cancel in their sum.
Thus the first diagonal, off-diagonal, and second diagonal entries of
$\mathsf S/|v|^{2K}$ have respective orders $O(1)$, $O(v^{-1})$,
$O(v^{-2})$.  This proves the degree bound, independently of any
cancellation at a central parameter value.  Its first diagonal is
positive for sufficiently large $|v|$.

The exact coefficient calculation in Appendix~\ref{cert:boundary}
proves $\det\mathsf S(v)>0$ for all real $v$ and all parameters in
$\mathcal D$.  A continuous Hermitian matrix with strictly positive
determinant cannot change its signature.  Its positive first diagonal
at large $|v|$ therefore implies $\mathsf S(v)>0$ on the whole line.
This completes the left-edge exclusion and the proof.
\end{proof}

\begin{lemma}[Winding count]\label{free:winding}
For $0\le\ell\le3$ the zero counts in Proposition~\ref{free:main} are
$2,1,0,0$, respectively.
\end{lemma}

\begin{proof}
Lemma~\ref{free:boundary} and the argument principle preserve total
zero multiplicity throughout $\rho:1\to0$.  At $\rho=1$,
Lemma~\ref{free:columns} identifies this count with the true matching
determinant by nonzero holomorphic scalar factors.  Individual
homotoped columns may vanish at interior singular factors of $P_K$;
these are zeros of a holomorphic determinant and do not affect the
boundary argument.  No assertion of interior nonvanishing is needed.

We may also move $(b,Z)$ to $(B/H_0,S/H_0)$, where
$H_0=10^8$, $B=33477607$, $S=270506819$.  This point lies in
$\mathcal D$, and the preceding exclusions hold on the connecting
segment.  At zero tail ratio, put $\lambda=-1/32+X$ and let
$(\widetilde x,\widetilde y)^t=P_K(0,1)^t$ in channel $h=1$.
The real polynomial
\[
 P_\ell(X)=\operatorname{Im}_{\mathrm{coeff}}
 (\widetilde x(X)\overline{\widetilde y}(X))
\]
is a nonzero constant multiple of the determinant; the bar here
conjugates coefficients.  Induction in the backward recurrence shows
that its degree is $15=2K-1$, with unscaled leading coefficient $-KZ$:
$\widetilde y$ has leading term $X^K$ and $\widetilde x$ has leading
term $-KiZ X^{K-1}$.  Write
\begin{equation}\label{free:axis-polynomials}
 P_\ell(iv)=\mathcal R_\ell(v^2)+iv\mathcal I_\ell(v^2).
\end{equation}
Both real polynomials have degree at most seven.

The rational signs in Appendix~\ref{cert:winding} give the following
root counts; all roots in the table are simple:
\[
\begin{array}{c|cc|cc}
 &\multicolumn{2}{c|}{\text{positive roots}}&
   \multicolumn{2}{c}{\text{negative roots}}\\
 \ell&\mathcal R_\ell&\mathcal I_\ell&\mathcal R_\ell&\mathcal I_\ell\\\hline
 0&5&6&2&1\\
 1&6&6&1&1\\
 2,3&7&7&0&0
\end{array}
\]
Here is why signs alone suffice.  Successive positive separators have
alternating one-coordinate sign changes, following the quadrant cycle
$++,-+,--,+-$.  They therefore give the listed distinct positive
roots, interlaced between the two polynomials.  For $\ell=0,1$ the
additional values at $-1/4,-4$ supply the listed negative roots.
These distinct roots exhaust each degree bound of seven.  Hence there
are no further roots, none repeated, and no unrecorded axis crossings.

The curve $P_\ell(iv)$ for $v>0$ crosses each successive axis in the
positive rotational direction and finishes in the first quadrant.
Let $n=5,6,7$ be the positive root count of $\mathcal R_\ell$.
Starting with the real value at zero and using the final odd-degree
asymptotic, the continuous argument change on $[0,\infty)$ is
$(n+1/2)\pi$.  This also follows by counting the consecutive quadrants
from the initial signs $-+$, $++$, $--$ in the three rows.  Conjugation
doubles that change on the full oriented imaginary axis.  A single
factor $iv-z$ contributes $+\pi$ if $\operatorname{Re}z<0$ and
$-\pi$ if $\operatorname{Re}z>0$.  Thus
\[
 N_{\mathrm{left}}-N_{\mathrm{right}}=2n+1,
 \qquad N_{\mathrm{left}}+N_{\mathrm{right}}=15,
\]
and $N_{\mathrm{right}}=(15-(2n+1))/2=2,1,0$.
The counting boundary has no zeros, so these are also the counts in
the closed half-plane.  The count transfers back by the homotopies.
\end{proof}

\subsection{Profile degree and nonvanishing of the exterior}

\begin{lemma}[Profile matching degree]\label{free:profile-degree}
The first conclusion of Proposition~\ref{free:main} holds.  More
precisely there is an invertible real-linear map
$\mathcal L:\mathbb R^2\to\mathbb C$ with negative determinant and
minimum stretch greater than $.048$ such that
\begin{equation}\label{free:profile-linearization}
 |j(b,Z)-\mathcal L(b-b_*,Z-Z_*)|<2.5\cdot10^{-11}
 \qquad((b,Z)\in\mathcal D).
\end{equation}
In particular $|j|>4.55\cdot10^{-10}$ on $\partial\mathcal D$.
\end{lemma}

\begin{proof}
Use $h=1$, $\ell=\lambda=0$, $K=34$ in
\eqref{free:backward}, dividing its $n$th factor by $n+1$, and call
the resulting product $L(b,Z)$.  Set
\[
 U_0=L(b_*,Z_*),\quad X_0=\partial_bL(b_*,Z_*),\quad
 Y_0=\partial_ZL(b_*,Z_*),\quad m_0=(U_0)_{01}.
\]
All matrix indices here are zero-based.  The exact rational
enclosures of Appendix~\ref{cert:profile} give
\begin{gather}
 9<|m_0|<10,\qquad
 (U_0)_{00}/m_0\in(2.835,2.837)+i(1.426,1.428),\notag\\
 |(U_0)_{11}/m_0|<10^{-12},\qquad
 (X_0)_{11}/m_0\in(-.0811,-.0810)+i(.1361,.1363),\notag\\
 (Y_0)_{11}/m_0\in(-.0001,.0001)+i(.1237,.1240),\notag\\
 \max_{i,j}|(X_0)_{ij}/m_0|,
 \max_{i,j}|(Y_0)_{ij}/m_0|<20,\notag\\
 \max_{i,j}|L_{ij}-(U_0)_{ij}-(b-b_*)(X_0)_{ij}
                   -(Z-Z_*)(Y_0)_{ij}|<50000|m_0|\delta^2,\notag\\
 |\det L|<6.2\cdot10^{-12}|m_0|^2.\label{free:profile-enclosures}
\end{gather}
The same certificate, propagating these estimates by rational
inequalities, gives
\begin{equation}\label{free:denominator-gap}
 |L_{01}-iZL_{00}/5|-Z|L_{00}|/5>.6|m_0|,
 \qquad |L_{01}/m_0-1|<5\cdot10^{-7}.
\end{equation}
Since the tail ratio lies in the disk with center $-iZ/5$ and radius
$Z/5$, the first component $L_{00}R_++L_{01}$ cannot vanish.
Lemma~\ref{free:columns} proves the desired nonvanishing of $H$ before
we take its logarithmic derivative.

The second-to-first component ratio is $j(b,Z)$.  Subtracting its value
at zero tail gives the exact identity
\[
 j-\frac{L_{11}}{L_{01}}
 =-\frac{R_+\det L}{L_{01}(L_{00}R_++L_{01})}.
\]
Using $|R_+|\le2Z/5$ and \eqref{free:denominator-gap} bounds its
modulus by
\[
 \frac{(2Z/5)|\det L|}{.6|m_0|\,|L_{01}|}.
\]
Define
\[
 \mathcal L(\eta,\zeta)=
 \eta\frac{(X_0)_{11}}{m_0}+\zeta\frac{(Y_0)_{11}}{m_0}.
\]
Appendix~\ref{cert:profile} gives the exact propagation of the displayed
bounds to \eqref{free:profile-linearization}, the minimum stretch
bound, and $|\mathcal L(\eta,\zeta)|\le.283\delta$ on the disk.
Its determinant is negative already from the two enclosing
rectangles.  The boundary size of the linear map exceeds
$.048\delta=4.8\cdot10^{-10}$; subtracting the error proves the
stated boundary margin.  Straight homotopy to $\mathcal L$ avoids
zero on the boundary, so its degree is $-1$.  Finally
$.283\delta+2.5\cdot10^{-11}<3\cdot10^{-9}$ proves the disk bound.
\end{proof}

\begin{lemma}[Exterior inequalities]\label{free:exterior}
For $(b,Z)\in\mathcal D$, $H(-ib,6,-iz)$ is nonzero for all $z\ge Z$,
and the inequalities in \eqref{free:exterior-bounds} and
$\operatorname{Im}j(b,z)>0$ for $z\ge3$ hold.  In radial variables,
for any nonzero constant multiple $Q(r)$ of $H(-ib,6,-ir^2/4)$,
\begin{equation}\label{free:radial-quotients}
 \frac{A'}A=-\frac r2\operatorname{Im}j(b,r^2/4),\qquad
 \frac wr=\frac12+\operatorname{Re}j(b,r^2/4),
 \quad A=|Q|,\quad w=r/2+2\partial_r\arg Q.
\end{equation}
\end{lemma}

\begin{proof}
Use the forward matrix \eqref{free:forward-matrix} at
$\lambda=\ell=0$, $s=iz$, $K=5$, and let
\[
 \mathsf H=T^*\begin{pmatrix}5&iz\\-iz&0\end{pmatrix}T
   =\begin{pmatrix}A_0&\overline c\\c&d_0\end{pmatrix}.
\]
The exact coefficient checks in Appendix~\ref{cert:exterior} prove
\begin{align}\label{free:exterior-forms}
 &d_0>0,\qquad .043d_0+\operatorname{Re}c-1130z>0
 &&(z\ge2.704),\notag\\
 &-\operatorname{Im}c-1130z>0&&(z\ge3),\notag\\
 &|\det T|=\prod_{n=0}^4|t_nd_n|<1130.
\end{align}
Every $b$ in the disk is positive, so all forward divisors are
nonzero.  The exact jet has nonpositive terminal cone form, hence
$J^*\mathsf HJ\le0$.  If its first component were zero, positivity
of $d_0$ would force its second component to be zero, contradicting
Lemma~\ref{free:columns}.  Thus $j$ is well defined.  Moreover
\[
 \det\mathsf H=-z^2|\det T|^2,
 \qquad \left|j+\frac c{d_0}\right|
 \le\frac{z|\det T|}{d_0}.
\]
Taking the largest possible real part and smallest possible imaginary
part in this disk, and using \eqref{free:exterior-forms}, proves the
assertions.  All $Z$ in \eqref{free:disk} exceed $2.704$.
Finally $Q'/Q=(ir/2)j$ gives \eqref{free:radial-quotients}.
\end{proof}

Proposition~\ref{free:main} now follows from
Lemmas~\ref{free:high-angular}, \ref{free:winding},
\ref{free:profile-degree}, and \ref{free:exterior}.
